\section{Reading What the Composer Says}
\label{sec:ch09-reading-the-errors}
\index{composition!reading the errors|(}%

Satisfiability is the check that worries people. It is not the check they meet.

\index{composition!error catalogue size}%
I counted the composition errors the library can produce, from the type of every
declaration in its error module at version 0.63.2. There are 123: 119 factory
functions that return an \code{Error}, and four exported constants that are one.
Exactly one of the 123 is satisfiability.

Counting them is fiddlier than it looks, and the fiddly part is the finding. The
module exports 203 things in total, and the other 80 are strings rather than
errors: message fragments, spliced into the errors above to build a sentence. So
counting declarations instead of their types gives 203, which is wrong by a
factor of nearly two. I know because that is the number I had first.

That number is not an invitation to read them. It is there to set an
expectation: a composer is mostly a very thorough schema linter that happens to
be reading two schemas at once, and the exotic distributed-systems check is one
line item in it.

The rest of this section is five failures, all produced from Mosaic's real pair
of schemas with one edit each, because an error message printed from an invented
fixture teaches you the fixture.

\subsection*{A field two subgraphs both claim}

\index{shareable@\code{@shareable}!the error it produces}%
Give Mosaic's \code{Product} a \code{title} and let both services publish one:

\begin{minted}[fontsize=\footnotesize,breaklines]{text}
The Object "Product" defines the same fields in multiple subgraphs without the "@shareable" directive:
 The field "title" is defined in the following subgraphs: "catalog", "mosaic".
 However, it is not declared "@shareable" in any of them.
\end{minted}

This is the rule
chapter~\ref{ch:the-federation-model} covered and
chapter~\ref{ch:the-first-cut-extracting-catalog} met three times before its tag
would ship. Two subgraphs may not both define a field unless one of them says it
is shareable~\autocite{apollo2026compositionrules}. The message names the type,
the field, both subgraphs and the missing directive, and there is nothing to
work out.

Hold on to its shape, because the next two errors have it and mean something
else.

\subsection*{A shared field with two different types}

Mark \code{averageRating} shareable in both services and give Catalog's an
\code{Int} where Mosaic has a \code{Float}:

\begin{minted}[fontsize=\footnotesize,breaklines]{text}
Each instance of a shared field must resolve identically across subgraphs.
The field "Product.averageRating" could not be federated due to incompatible types across subgraphs.
The discrepancies are as follows:
 The named type "Int" is returned by the following subgraph: "catalog".
 The named type "Float" is returned by the following subgraph: "mosaic".
\end{minted}

Apollo's second composition rule: a shared field needs a compatible return type
and compatible argument types everywhere it is
defined~\autocite{apollo2026compositionrules}. Also clear, also nothing to work
out.

\index{composition!duplicated error messages}%
One thing about it is worth a sentence, because it will otherwise cost you a
minute of counting. The composer printed that error twice, identically, for one
field with one problem. So the number of blocks in the table is not the number
of things wrong with your graph.

\subsection*{An enum that drifted, and the only error that tells you what to do}

\index{composition!enum drift}%
Declare \code{ProductCategory} in both subgraphs with different members. Catalog
uses it as an output on \code{Product.category} and as an input in its filter,
which is the condition that makes this strict:

\begin{minted}[fontsize=\footnotesize,breaklines]{text}
Enum "ProductCategory" was used as both an input and output but was inconsistently defined across inclusive subgraphs. To update an Enum used as both an input and output, add any new Enum values with the @inaccessible directive in the origin subgraph. Next, add those new Enum values to all other subgraphs that define the Enum—this time without the @inaccessible directive. Finally, once all subgraphs have been updated, remove @inaccessible from the Enum values in the origin subgraph.
\end{minted}

\index{composition!errors that prescribe a migration}%
That is the only message in this chapter that tells you what to do rather than
what is wrong, and it is prescribing a three-step migration across two
deployments because there is no safe single step. An enum read as output wants
to grow; an enum accepted as input wants to shrink; and a value that exists on
one side and not the other is unsafe in one direction or the other whichever way
you add it.

\subsection*{The two that name the wrong thing}

\index{composition!errors that name the symptom}%
Here is where reading these gets a skill attached to it.

Take the \code{@key} off Mosaic's \code{Product} entirely, so that the type is
declared in both services and is an entity in only one:

\begin{minted}[fontsize=\footnotesize,breaklines]{text}
The Object "Product" defines the same fields in multiple subgraphs without the "@shareable" directive:
 The field "id" is defined and declared "@shareable" in the following subgraph: "catalog".
 However, it is not declared "@shareable" in the following subgraph: "mosaic".
\end{minted}

Nothing in that message mentions a key.

\index{key@\code{@key}!implicitly shareable}%
The reasoning is sound and it runs backwards from where you would look. A key
field is shareable by definition, because every subgraph declaring the entity
has to be able to answer with it, which is why
chapter~\ref{ch:the-first-cut-extracting-catalog} could put \code{id} on both
sides of the cut and compose. Take the key away and \code{id} stops being a key
field. It becomes an ordinary field that two subgraphs both declare, one of
which is still implicitly shareable and one of which is now nothing at all. So
the composer reports shareability, because shareability is the rule that broke
first.

Now key the two services on different fields. Mosaic keys on \code{sku} and
publishes one; Catalog still keys on \code{id}:

\begin{minted}[fontsize=\footnotesize,breaklines]{text}
The Object "Product" defines the same fields in multiple subgraphs without the "@shareable" directive:
 The field "id" is defined and declared "@shareable" in the following subgraph: "catalog".
 However, it is not declared "@shareable" in the following subgraph: "mosaic".
 The field "sku" is defined and declared "@shareable" in the following subgraph: "mosaic".
 However, it is not declared "@shareable" in the following subgraph: "catalog".
\end{minted}

Symmetric, and again never the words \enquote{your keys disagree}. Each side's
key field is implicitly shareable at home and an ordinary duplicated field
abroad, so the composer reports it twice, once in each direction.

\index{composition!three causes one error}%
Three of the five failures in this section come out of a single error factory:

\begin{minted}[fontsize=\footnotesize,breaklines]{javascript}
function invalidFieldShareabilityError(objectData, invalidFieldNames) { ... }
\end{minted}

Its two message shapes are visible in the source, one for a field shareable in
none of the subgraphs and one for a field shareable in some and not others.
Three causes, one error, and only the first of the three is really about
shareability.

\medskip
So the rule I would give somebody reading composition output for the first time
is this. \textbf{The composer tells you where it noticed, not what you did.} It
compares two schemas coordinate by coordinate and reports the first rule that
fails at each one, and the rule that fails is often downstream of the mistake.
An error about a field you never thought about, on a type you did not touch, is
the normal case rather than a sign that something strange is happening.

The practical form of that: when a shareability error names a key field, stop
reading it as a shareability error. It is telling you something is wrong with
the entity, and the entity's key is the first thing to look at.

All of which assumes you want the composer to refuse. There is a flag that
makes it stop.

\index{composition!reading the errors|)}%
